\documentclass[10pt]{article}
\usepackage[letterpaper,margin=0.85in]{geometry}
\input{glyphtounicode}
\pdfgentounicode=1
\pdftrailerid{}
\usepackage{amsmath,amssymb,amsthm}
\usepackage{booktabs}
\usepackage[colorlinks=true,linkcolor=blue,citecolor=blue,urlcolor=blue]{hyperref}
\hypersetup{pdftitle={An infinite family of paired Hadamard matrices from Williamson-type matrices},pdfauthor={Jeffery Kline}}
\newtheorem{theorem}{Theorem}
\newtheorem{lemma}[theorem]{Lemma}
\newcommand{\tr}{^{\mathsf T}}
\emergencystretch=2em

\title{An infinite family of paired Hadamard matrices\\from Williamson-type matrices}
\author{Jeffery Kline}
\date{November 27, 2026\\[3pt]\small Version 0.1.0, \href{https://doi.org/10.5281/zenodo.24615122}{doi:10.5281/zenodo.24615122}}

\begin{document}
\maketitle

\begin{abstract}
Every Williamson-type quadruple of positive odd order $w$ gives a Hadamard matrix of order $12w$ whose rows and columns can be paired into two-by-two Hadamard blocks. We give one pairing of the order-12 Cooper--Wallis design that works for every such quadruple. Turyn's quadruples give infinitely many distinct paired orders. When a conference matrix of order $6w+2$ exists, the published paired-block construction gives a Hadamard matrix of order $8(6w+1)$. At $w=331$ we construct the matrix of order 15,896 and verify every pair of rows. The theorem also gives orders 18,008, 18,296 and 23,096. We found no earlier construction at these four orders in the sources searched through November 27, 2026.
\end{abstract}

\section{Introduction}

A Hadamard matrix of order $N$ is an $N\times N$ matrix $H$ with entries $\pm1$ and $HH\tr=NI$. Call a Hadamard matrix \emph{paired-block} if its rows and its columns can each be partitioned into ordered pairs so that every $2\times2$ block on a row pair and a column pair is Hadamard. A conference matrix of order $c$ has zero diagonal, $\pm1$ off the diagonal, and row Gram matrix $(c-1)I$.

The pairing in Section~2 produces paired-block Hadamard matrices of infinitely many distinct orders. For example, take $q=5^k$ with $k\ge1$. Each $q$ is a prime power congruent to $1$ modulo $4$, so Turyn's construction~\cite{Tu72} supplies a Williamson-type quadruple of positive odd order $w=(q+1)/2$. The resulting paired orders $12w$ increase with $k$. Passing from these paired orders to orders $8(6w+1)$ requires a conference matrix for the chosen $w$.

Theorem~1 of the order-4792 paper~\cite{Kl4792} gives a Hadamard matrix of order $8m$ from a symmetric or skew-symmetric conference matrix of order $m+1$ and a paired-block Hadamard matrix of order $2(m-1)$, for odd $m\ge3$. That paper supplies pairs for the order-52 Cooper--Wallis design when two of the four Williamson-type members are equal. We give pairs for the order-12 design that work for every Williamson-type quadruple. Taking $m=6w+1$ for a quadruple of order $w$ then gives Hadamard matrices of orders 15,896, 18,008, 18,296 and 23,096. The release includes the matrix of order 15,896 and an all-pairs verification of its Gram matrix. We found no earlier construction of the four orders in the sources searched through November 27, 2026.

\section{The order-12 pairing}

Let $A,B,C,D$ be the circulants with first rows $(1,0,0),(0,1,0),(0,0,1),(0,0,0)$, respectively. If a first row is $a$, its circulant has entry $a_{j-i\bmod 3}$ at $(i,j)$ for indices $0,1,2$. Use the back-diagonal reversal $R_{ij}=1$ for $i+j=2$ and $R_{ij}=0$ otherwise. In the Goethals--Seidel layout~\cite{GS70} put
\[
G(A,B,C,D)=
\begin{pmatrix}
 A & BR & CR & DR\\
 -BR & A & -D\tr R & C\tr R\\
 -CR & D\tr R & A & -B\tr R\\
 -DR & -C\tr R & B\tr R & A
\end{pmatrix}.
\]
Form $E_1,E_2,E_3,E_4$ by applying $G$ to the four ordered inputs
\[
(A,B,C,D),\qquad(B,-A,D,-C),\qquad(C,-D,-A,B),\qquad(D,C,-B,-A).
\]
The four first rows above are the length-three entry in Cooper and Wallis's table following Corollary~4~\cite{CW72}. For an implementation of the Cooper--Wallis array, see SageMath~\cite{Sage}. The supports of the $E_i$ partition the entries, with nonzero coefficients $\pm1$. Indeed, each $3\times3$ block uses the three disjoint circulant supports with the same transpose and reversal operations. The Cooper--Wallis identity, read as the Gram of $L=\sum_{i=1}^4 x_i E_i$ in commuting indeterminates, is $LL\tr=3(\sum_i x_i^2)I_{12}$. Here $A,B,C$ are permutation matrices and $D=0$, giving $AA\tr+BB\tr+CC\tr+DD\tr=3I_3$. Comparing coefficients gives
\begin{equation}\label{eq:design}
E_iE_i\tr=3I_{12},\qquad E_iE_j\tr+E_jE_i\tr=0\quad (i\ne j).
\end{equation}

\begin{lemma}\label{lem:pairing}
Use the following ordered pairs of zero-based indices of $L$, numbering each axis consecutively across the blocks of size three.
\[
\begin{array}{@{}lrrrrrr@{}}
\text{Rows} &(0,4)&(1,10)&(2,7)&(3,5)&(6,8)&(9,11)\\
\text{Columns} &(1,10)&(3,2)&(4,7)&(6,9)&(8,0)&(11,5)
\end{array}
\]
For every scalar sign assignment $x_i\in\{1,-1\}$, all 36 blocks on these pairs are Hadamard.
\end{lemma}

\begin{proof}
Write an entry $\pm x_i$ as $\pm i$. Reading the four entries of each block in row-major order gives the following table, with rows and columns in the displayed pair orders.
{\small
\[
\begin{array}{@{}rrrrrr@{}}
(2,3,2,-3)&(4,3,3,-4)&(1,-4,1,4)&(1,-2,2,1)&(-3,1,-3,-1)&(-4,-2,-4,2)\\
(1,-4,4,1)&(1,2,-1,2)&(-2,-3,3,-2)&(-4,3,4,3)&(1,3,1,-3)&(-2,4,2,4)\\
(3,-2,3,2)&(-2,1,-2,-1)&(4,1,-4,1)&(-3,-4,3,-4)&(-4,2,2,4)&(3,1,-1,3)\\
(-1,1,-4,-4)&(1,2,2,-1)&(2,2,3,-3)&(-3,-4,4,-3)&(4,-4,2,2)&(-3,3,1,1)\\
(4,-4,-1,-1)&(3,3,-4,4)&(-2,2,3,3)&(1,-1,2,2)&(3,-1,1,3)&(2,-4,-4,-2)\\
(-3,2,2,3)&(4,4,3,-3)&(-1,4,4,1)&(1,1,-2,2)&(-2,2,4,4)&(3,3,1,-1)
\end{array}
\]}
Each tuple has an odd number of minus signs, and each variable index has even multiplicity. Thus $abcd=-1$ for all 16 assignments, where $a,b,c,d$ are the entries in row-major order. This gives $ac+bd=0$.
\end{proof}

\begin{theorem}\label{thm:general}
Let $w$ be a positive odd integer with a Williamson-type quadruple of order $w$. If a real conference matrix of order $6w+2$ exists, then a real Hadamard matrix of order $8(6w+1)$ exists.
\end{theorem}

\begin{proof}
Write the quadruple as $W_1,\dots,W_4$, sign matrices of order $w$ with $W_iW_j\tr=W_jW_i\tr$ for all $i,j$ and $\sum_i W_iW_i\tr=4wI_w$. Put $F=\sum_i E_i\otimes W_i$. For each outer row pair $(r,s)$ and each inner index $u$, pair $(r,u)$ with $(s,u)$. For each outer column pair $(c,d)$ and inner index $v$, pair $(c,v)$ with $(d,v)$. The four entries at these pairs evaluate the corresponding block of $L$ at $x_i=(W_i)_{uv}$. Lemma~\ref{lem:pairing} gives a Hadamard block at every pair. In particular the pairing works for noncirculant members.

In $FF\tr$, the terms indexed by $(i,j)$ and $(j,i)$ cancel for $i\ne j$ by amicability and~\eqref{eq:design}. The diagonal terms sum to $3I_{12}\otimes\sum_i W_iW_i\tr=12wI_{12w}$. Thus $F$ is paired-block of order $12w$. Put $m=6w+1$, so $F$ has order $2(m-1)$ and the conference has order $m+1$. By Delsarte, Goethals and Seidel~\cite[Corollary~2.2]{DGS71}, the conference matrix is equivalent to one of symmetric or skew type. Theorem~1 of~\cite{Kl4792} applies and gives order $8m$.
\end{proof}

\section{Four orders}

For a prime power $m=6w+1$, odd $w$ makes $m\equiv3\pmod4$, so Paley's construction~\cite{Pa33} gives the requisite skew conference matrix. For a prime power $q\equiv1\pmod4$, Turyn~\cite{Tu72} gives a quadruple of order $(q+1)/2$ and Spence~\cite[Theorem~4]{Sp77} gives one of order $q^r(q+1)/2$ for $r\ge0$. The parameters below satisfy the two prime-power conditions separately. The entries in the $m$ column are all prime.
\begin{center}
\begin{tabular}{@{}rrrrl@{}}
\toprule
Hadamard order $8m$ & $m$ & Conference order $m+1$ & $w$ & Williamson-type source\\
\midrule
15,896 & 1987 & 1988 & 331 & Turyn, $q=661$\\
18,008 & 2251 & 2252 & 375 & Spence, $q=5$, $r=3$\\
18,296 & 2287 & 2288 & 381 & Turyn, $q=761$\\
23,096 & 2887 & 2888 & 481 & Turyn, $q=31^2$\\
\bottomrule
\end{tabular}
\end{center}

\section{The matrix of order 15,896}

The construction uses Turyn's quadruple of order 331, the paired matrix of order $12\cdot331=3972$, and the Paley skew conference of order $1988$. For the latter, set $Q_{xy}=\chi(y-x)$ for $x,y\in\mathbb F_{1987}$, with $\chi$ the quadratic character extended by $\chi(0)=0$. Border $Q$ by a top row of ones, a left column of minus ones and a zero corner. Applying the order-4792 construction to these inputs gives $8\cdot1987=15896$. The program \texttt{evidence/build.cpp} makes the Turyn input over $\mathbb F_{661}[x]/(x^2-2)$ with primitive element $2+x$, checks the Williamson-type identities, and writes the packed matrix. Appendix~\ref{app:recipe} specifies the same matrix entry by entry.

The uncompressed payload of \texttt{evidence/H15896.packed.gz} has $31{,}585{,}352$ bytes. It holds $15{,}896$ rows of $1{,}987$ bytes. Entries are packed in row-major order, most significant bit first, with $0$ for $+1$ and $1$ for $-1$.

The program \texttt{evidence/verify.cpp} reads only this payload and checks its exact length and the Hamming distance of every two distinct rows. For sign rows the inner product is $N-2d$, where $d$ is their Hamming distance. All $\binom{15896}{2}=126{,}333{,}460$ unordered row pairs have distance $7{,}948$. The diagonal inner products are $15{,}896$ because each entry is a sign. This verifies $HH\tr=15896I$ directly from the bytes.

\section{Earlier orders and scope of the search}

Miyamoto's Theorem~7~\cite{Mi91} takes a conference matrix of order $m+1$ and a quaternary Hadamard matrix of order $m-1$ to give a quaternary Hadamard matrix of order $4m$, which expands to a real Hadamard matrix of order $8m$. A quaternary Hadamard matrix $Y$ of order $k$ has entries in $\{\pm1,\pm i\}$ and satisfies $YY^*=kI$. Zaitsev~\cite{Zai} publishes quaternary matrices of orders 1566 and 1782. With the Paley conferences for $m=1567$ and $m=1783$, these give the earlier orders 12,536 and 14,264. Zaitsev lists the quaternary order 1986 as open in his table, and the three quaternary orders 2250, 2286 and 2886 lie beyond its range.

Kumar's pinned update~\cite{Ku} to Cati and Pasechnik's Table~4~\cite{CP} lists exponents $4,4,4,5$ at the four odd parts $1987,2251,2287,2887$, respectively, in its table of orders $2^a m$. A later release supplies $16\cdot2887$~\cite{KlLarge}. The construction here gives exponent $3$ at all four odd parts. The 2026 survey's list of real orders~\cite{EKLO} runs through $10{,}000$. The release repository records the dated searches of these tables, public repositories and literature indices.

Below 10,000, the input families tabulated in the supplement give 79 distinct final orders $8(6w+1)$, all previously known. Of these, 27 already follow from Turyn or Spence with a Paley conference, and the scalar input gives order 56. A recent overlap is $5912$ ($w=123$), which Kumar constructed by Spence's method~\cite{Ku}. His main-branch certificate is dated October~29, 2026, four weeks before this note. Order $8696$ ($w=181$) also appears on the November survey's open list, although an older Williamson-type construction covers it~\cite{Ku,EKLO}. For odd $w$ every final order is $56\bmod 96$. This excludes 428 and all 29 entries in the survey's Table~4 row for sporadic computer searches, each of which is $4\bmod 8$. The paired orders $12w$ give a different comparison: they include $2076$ and $9708$, both discussed in Section~4 of the survey~\cite{EKLO}.

\section{Role of AI}
AI assisted the derivations, the computer verification, the literature search and the exposition under the author's direction. Process-separated AI reviews checked the proofs and certificates; they are not independent expert peer review, and no independent human expert has examined the proofs. The release repository records the sources, the reviews and their dispositions.

\appendix
\section{Reconstructing the order-15,896 matrix}\label{app:recipe}

The following choices specify each entry of the certificate. All indices start at zero. Write $B=\left(\begin{smallmatrix}1&1\\1&-1\end{smallmatrix}\right)$ and let $\chi_p$ be the quadratic character on $\mathbb F_p$, with $\chi_p(0)=0$.

\subsection*{Four circulant rows and their paired tiles}

Set $q=661$, $w=331$ and work in $\mathbb F_q[\alpha]/(\alpha^2-2)$ with $g=2+\alpha$. For $0\le j\le661$, put
\[
 b_j=\chi_q\!\left(( (g^j)^q-g^j )/g^{331}\right).
\]
The quotient lies in the base field. For $0\le k<w$, put $r_k=(-1)^k b_{2k}$ and $o_k=(-1)^k b_{2k+1}$. Set $s_k=o_{(k+165)\bmod w}$ and let $e=(1,0,\ldots,0)$ of length $w$. The first rows of $W_1,W_2,W_3,W_4$, in that order, are $s,s,r+e,r-e$. For each matrix, entry $(u,v)$ is its first-row entry at $(v-u)\bmod w$.

Set $m=1987$, $n=(m-1)/2=993$, and $d=2n=6w=1986$. For $0\le x,y<d$, write $x=aw+u$ and $y=bw+v$ with $0\le u,v<w$. Take the tuple in row $a$ and column $b$ of the $6\times6$ tuple table in the proof of Lemma~\ref{lem:pairing}. Replace each signed label $\pm i$ by $\pm(W_i)_{uv}$ and read the result in row-major order as a $2\times2$ matrix $T_{xy}$. Define the second array of tiles by
\[
 U_{xy}=-\tfrac12 (B T_{yx} B)^{\mathsf T}.
\]
Both arrays consist of signs. The exchanged tile indices and the transpose within the tile both enter this formula.

\subsection*{Conference core and borders}

List the integers from $1$ to $1986$ in two groups, the quadratic residues modulo $m$ first, then the nonresidues, each group in increasing order. Denote the list by $p_0,\ldots,p_{d-1}$. Put $C_{xy}=\chi_m(p_y-p_x)$ and form a $2\times2$ matrix of sums
\[
 M_{sa}=\sum_{j=0}^{n-1} C_{sn,\,an+j}\qquad(s,a\in\{0,1\}).
\]
The first row of each half suffices because multiplication by a quadratic residue permutes each half transitively and preserves $C$.

On each axis the first eight indices are the border, numbered $0,\ldots,7$. The remaining indices are the tuples $(s,h,c,t,z)$ in lexicographic order, with $s,h,c,z\in\{0,1\}$ and $0\le t<n$. Thus $z$ varies fastest. A row tuple has group index $i=4s+2h+c$ and point index $x=sn+t$. For a column tuple $(a,k,b,u,v)$ the corresponding indices are $j=4a+2k+b$ and $y=an+u$.

Define three $8\times8$ matrices, using these binary decompositions for $i$ and $j$ and writing $[P]$ for the indicator of $P$:
\[
 \begin{aligned}
 R_{ij} &= B_{s k}\bigl(1-2[\,2h+c=2a+b\,]\bigr),\\
 K_{ij} &= [h=k]\,B_{sa}B_{cb},\\
 S_{ij} &= [c=b]\,M_{sa}B_{hk}.
 \end{aligned}
\]
Use ordinary matrix multiplication. The upper left corner is $RK/2$. The entry from border row $\ell$ to an interior column in group $j$ is $R_{\ell j}$. The entry from an interior row in group $i$ to border column $\ell$ is $-(SK)_{i\ell}$.

\subsection*{Interior entries and serialization}

For an interior row and column with the tuple coordinates above, the entry is
\[
 \begin{cases}
 C_{xy} B_{hk}, & c=b,\ x\ne y,\\
 (-1)^{z+v}B_{hk}, & c=b,\ x=y,\\
 (-1)^{z+v}(T_{xy})_{hk}, & c=0,\ b=1,\\
 (-1)^{z+v}(U_{xy})_{hk}, & c=1,\ b=0 .
 \end{cases}
\]
There are $8+16n=15{,}896$ rows and columns. Traverse columns in the stated order within each row and rows in the same order. Encode $+1$ as zero and $-1$ as one, with the first bit in each byte the most significant. Each row fills exactly 1,987 bytes. The headerless, uncompressed stream has 31,585,352 bytes and SHA-256
\begin{center}\small\ttfamily
93ba6e82b4e8a137a00b2e6898fa03466cd1cb5eb82bca245cefde5eafc77dcf
\end{center}
The digest identifies the precise representative and its ordering. An independent orthogonality check requires that every two distinct rows differ in exactly 7,948 positions.

\begin{thebibliography}{99}\small
\bibitem{CP} M. Cati and D. V. Pasechnik, A database of constructions of Hadamard matrices, \href{https://arxiv.org/abs/2411.18897v2}{arXiv:2411.18897v2} (2025).
\bibitem{CW72} Joan Cooper and Jennifer Wallis, A construction for Hadamard arrays, \emph{Bull. Austral. Math. Soc.} \textbf{7} (1972), 269--277, \href{https://doi.org/10.1017/S0004972700045019}{doi:10.1017/S0004972700045019}.
\bibitem{DGS71} P. Delsarte, J. M. Goethals and J. J. Seidel, Orthogonal matrices with zero diagonal. II, \emph{Canad. J. Math.} \textbf{23} (1971), 816--832, \href{https://doi.org/10.4153/CJM-1971-091-x}{doi:10.4153/CJM-1971-091-x}.
\bibitem{EKLO} R. Egan, I. S. Kotsireas, S. London and P. \'O~Cath\'ain, Recent progress in the construction of Hadamard matrices, \href{https://arxiv.org/abs/2610.27175v2}{arXiv:2610.27175v2} (2026).
\bibitem{GS70} J. M. Goethals and J. J. Seidel, A skew Hadamard matrix of order 36, \emph{J. Austral. Math. Soc.} \textbf{11} (1970), 343--344.
\bibitem{Kl4792} J. Kline, A Hadamard matrix of order 4792, version 0.1.0 (2026), \href{https://doi.org/10.5281/zenodo.24597897}{doi:10.5281/zenodo.24597897}.
\bibitem{KlLarge} J. Kline, New Hadamard orders above $10^4$, version 0.1.0 (2026), \href{https://doi.org/10.5281/zenodo.24517277}{doi:10.5281/zenodo.24517277}.
\bibitem{Ku} K. Kumar, \href{https://github.com/vatskishan03/Hadamard/blob/d885022144ba6bf4efa8315e14871a3736947527/FRONTIER.md\#L58-L62}{Updated Table 4}, Hadamard repository, commit \texttt{d885022} (accessed November 27, 2026).
\bibitem{Mi91} M. Miyamoto, A construction of Hadamard matrices, \emph{J. Combin. Theory Ser. A} \textbf{57} (1991), 86--108, \href{https://doi.org/10.1016/0097-3165(91)90008-5}{doi:10.1016/0097-3165(91)90008-5}.
\bibitem{Pa33} R. E. A. C. Paley, On orthogonal matrices, \emph{J. Math. Phys.} \textbf{12} (1933), 311--320.
\bibitem{Sage} The Sage Developers, SageMath, modules \href{https://doc.sagemath.org/html/en/reference/combinat/sage/combinat/matrices/hadamard_matrix.html}{Hadamard matrices} and \href{https://doc.sagemath.org/html/en/reference/combinat/sage/combinat/t_sequences.html}{T-sequences}.
\bibitem{Sp77} E. Spence, An infinite family of Williamson matrices, \emph{J. Austral. Math. Soc. Ser. A} \textbf{24} (1977), 252--256, \href{https://doi.org/10.1017/S1446788700020255}{doi:10.1017/S1446788700020255}.
\bibitem{Tu72} R. J. Turyn, An infinite class of Williamson matrices, \emph{J. Combin. Theory Ser. A} \textbf{12} (1972), 319--321, \href{https://doi.org/10.1016/0097-3165(72)90095-7}{doi:10.1016/0097-3165(72)90095-7}.
\bibitem{Zai} V. Zaitsev, \href{https://www.cs.uleth.ca/~zaitsevv/quaternary.html}{Quaternary Hadamard matrices}, online table (accessed November 27, 2026).
\end{thebibliography}
\end{document}
